\section{An agentic review system}\label{sec:system}

\begin{figure}[t]\centering
\resizebox{\linewidth}{!}{% Fig: the agentic review system. Included from sections/system.tex.
\begin{tikzpicture}[
  font=\sffamily\scriptsize, node distance=4mm and 3.2mm,
  box/.style={draw=black!45, rounded corners=2pt, fill=#1, align=center, inner sep=3pt, minimum height=11mm, text width=19mm},
  box/.default=white,
  src/.style={box=blue!6}, agent/.style={box=green!8}, check/.style={box=orange!10}, upbox/.style={box=red!6},
  arr/.style={-{Stealth[length=2mm]}, draw=black!55, thick}]
\node[src] (fetch) {\textbf{Fetch}\\UniProt record, GOA rows, cached papers};
\node[src, right=of fetch] (dr) {\textbf{Deep research}\\several providers, cited reports};
\node[agent, right=of dr] (rev) {\textbf{Agent review}\\action + reason + verbatim quotes per row};
\node[check, right=of rev] (val) {\textbf{Validate}\\schema, term ids, quotes, policies};
\node[check, right=of val] (pr) {\textbf{Pull request}\\review, merge, history record};
\node[upbox, below=6mm of val, text width=40mm, minimum height=9mm] (up) {\textbf{Upstream}: source defects, NOT candidates,\\proposed terms, mapping and rule edits};
\node[agent, below=6mm of dr, text width=40mm, minimum height=9mm] (tools) {\textbf{Tools}: OLS, PubMed, UniProt, PANTHER/PAINT,\\InterPro, reproducible bioinformatics};
\draw[arr] (fetch) -- (dr); \draw[arr] (dr) -- (rev); \draw[arr] (rev) -- (val); \draw[arr] (val) -- (pr);
\draw[arr] (val.north) to[bend right=25] node[above, font=\sffamily\tiny]{fails: revise} (rev.north);
\draw[arr] (pr.south) |- (up.east);
\draw[arr] (tools.east) -| (rev.south);
\end{tikzpicture}
}
\caption{The agentic review system. Each gene review passes fetch, deep research, agent review,
validation and pull-request review; validated findings feed reports to upstream resources.}
\label{fig:system}
\end{figure}

\subsection{The review document}
One YAML document per gene (LinkML schema, \texttt{gene\_review.yaml}): a standalone
description, the references used, one entry per existing annotation with an action, summary,
reason and \texttt{supported\_by} quotes, the gene's \emph{core functions} (MF + process +
location + complex), proposed new terms, knowledge gaps, suggested experiments.

\subsection{The action vocabulary}
\begin{description}[nosep,font=\normalfont\scshape]
  \item[accept] correct and part of the core function.
  \item[keep\_as\_non\_core] correct and as precise as warranted, but peripheral.
  \item[modify] sound but a better term exists (more specific, or sign/lateral correction);
        carries proposed replacements.
  \item[too\_general] true but entailed by a more specific term in the review.
  \item[mark\_as\_over\_annotated] not strictly false, but an over-reach.
  \item[remove] unlikely to be correct on the combined evidence.
  \item[undecided] evidence insufficient or inaccessible.
  \item[new] proposed annotation absent from GOA.
\end{description}
\note{Make explicit: actions judge \emph{claims}, not evidence codes; rows sharing a term must be
consistent; ``correct but general'' is \act{MODIFY} not \act{KEEP\_AS\_NON\_CORE}.}

\subsection{Guardrails against hallucination}
\begin{itemize}[nosep]
  \item Cached sources only; every \texttt{supporting\_text} must be a verbatim substring of the
        cached paper, UniProt record, Reactome entry or analysis file.
  \item Reference titles are checked against the fetched record (catches transposed PMIDs).
  \item Author-supplied term ids are validated against ontology branches; GOA ids are trusted.
  \item Policy validators encode curation conventions (protein binding, ND, root organelle
        terms, IBA reading, ``do not add what curators declined to add'').
  \item Experimental annotations whose full text is not cached are not removed (\act{UNDECIDED}).
  \item Bioinformatics analyses are reproducible scripts; results never hardcoded.
  \item Every change is a reviewed pull request with an append-only history record.
\end{itemize}

\subsection{Gene selection}
\todo{Describe honestly: not a random sample. Mix of (a) whole-proteome campaigns (e.g.\
\textit{P. putida}, \textit{S. pombe}, yeast), (b) thematic projects (proteostasis, immunity,
\ldots), (c) term-centric audits (all annotations to a term slated for obsoletion), (d) human
gene sets. Quantify selection; report results per campaign type to show robustness.}

\subsection{Models and agents}
\todo{Which agent frameworks and model versions, over what period; deep-research providers;
cost per gene; wall-clock.}
